% ============================================================
% 片段：s07  对应原文 appendix2.tex 第 1--20 行；PDF 第 68 页
% 内容：附录 B  群上同调
% 译者：lead   审校状态：初稿
% ============================================================

\section{群上同调}

\begin{keytakeaways}
  \begin{itemize}[leftmargin=1.4em,itemsep=2pt,topsep=2pt]
    \item \textbf{物理图像}：群上同调是“对称性在离散群上的反常与扩张”的统一语言：它把规范变换的闭合性/相容性表述为上闭链、把规范冗余表述为上边缘，从而把物理上的反常分类化为上同调群的计算。
    \item \textbf{数学工具}：$G$-模、$n$-上链群 $C^n(G,M)$、上边缘映射 $\delta_n$、上同调群 $H^n(G,M)$，以及与奇异上同调的同构 $H^n_{\rm grp}(G,A)\cong H^n_{\rm sing}(BG,A)$。
    \item \textbf{关键结论}：$H^3_{\rm grp}(\Pi_1,\Pi_2)$ 中的类就是正文第 5 节反复出现的 Postnikov 类 $\beta$ 的栖身之处；2-群的 Postnikov 数据 $(\Pi_1,\Pi_2,\alpha,\beta)$ 正是用它来刻画“$0$-形式对称性如何缠绕 $1$-形式对称性”。
  \end{itemize}
\end{keytakeaways}

\label{App:GrpChmlg}

群上同调是研究群及其作用的性质的数学工具。它是代数拓扑的一个分支，处理带有群作用的空间的上同调。
群系数在模中的上同调群定义如下。

设 $G$ 是一个群，$M$ 是一个 $G$-模。$G$ 以 $M$ 为系数的第 $n$ 个上同调群定义为商群
\[ H^n(G, M) = \frac{\text{ker}(\delta_n)}{\text{im}(\delta_{n-1})}, \]
其中 $\delta_n:H^n(G, M) \rightarrow H^{n+1}(G, M)$ 是由下式定义的上边缘映射：
\be\delta_n(f)(g_0, g_1, \ldots, g_n) = g_0f(g_1, \ldots, g_n) + \sum_{i=1}^{n}(-1)^if(g_0, \ldots, g_{i-1}g_i, \ldots, g_n) + (-1)^{n+1}f(g_0, \ldots, g_{n-1})\ee
这里 $f \in C^n(G, M)$，而 $C^n(G, M)$ 是 $G$ 以 $M$ 为系数的 $n$-上链群。

值得注意的是，有一个定理给出群上同调与奇异上同调之间的同构：
\be
H^n_{{\rm grp}}(G,A) \cong  H^n_{{\rm sing}}({\rm B}G,A) \ ,
\ee
其严格表述与证明可参见文献~\cite{weibel_1994}。

\begin{yizhu}
本附录的表述有一处不够严谨：$\delta_n$ 定义在 $n$-上链上，写作 $\delta_n:C^n(G,M)\to C^{n+1}(G,M)$，
而非原文所写的 $\delta_n:H^n\to H^{n+1}$（上边缘映射是上链层面的算子，$H^n$ 只是它诱导到商群上的映射）。
另请注意 $G$ 的作用是左作用，故上闭链条件中的 $g_0f(g_1,\dots)$ 表示用 $g_0$ 作用在系数上。
\end{yizhu}
